\documentclass[../../main.tex]{subfiles}
\begin{document}
\section{习题}\label{section:ch16-problems}

\begin{enumerate}
\item\label{exercise:ch16-01} 设 $V$ 是一个 $n$ 维向量空间。证明：$\det(v_1,\dots,v_n)$ 的绝对值至多差一个符号，就是由 $v_1,\dots,v_n$ 张成的平行多面体的 $n$ 维体积。
\item\label{exercise:ch16-02} 设 $M$ 是一个定向光滑 $n$ 流形，$\omega$ 是 $M$ 上一个紧支集光滑 $n$-形式。证明：$\int_M\omega$ 不依赖于用来定义它的定向坐标卡图册。
\item\label{exercise:ch16-03} 设 $F:M\to N$ 是一个定向光滑流形之间的光滑映射。证明：拉回 $F^*$ 保持闭形式与恰当形式。
\item\label{exercise:ch16-04} 设 $G$ 是一个紧 Lie 群。证明：其上的左不变 Haar 积分在差一个常数因子的意义下唯一。
\item\label{exercise:ch16-05} 用球坐标计算单位球面 $S^{n-1}\subseteq\mathbb{R}^n$ 的面积。
\item\label{exercise:ch16-06} 设 $D\subseteq\mathbb{R}^n$ 是一个积分区域，$G:D\to\mathbb{R}^n$ 是一个在边界上线性的映射。证明变量替换公式。
\item\label{exercise:ch16-07} 用 Stokes 定理计算正方形边界上一个 $1$-形式的积分。
\item\label{exercise:ch16-08} 通过在单位球面上积分一个适当的 $n-1$-形式，导出球面面积公式。
\item\label{exercise:ch16-09} 设 $\omega$ 是 $\mathbb{R}^2\setminus\{0\}$ 上的闭 $1$-形式。证明 $\omega$ 是恰当的当且仅当它绕单位圆周的积分为零。
\item\label{exercise:ch16-10} 证明：在一个单连通流形上，每个闭 $1$-形式都是恰当的。
\item\label{exercise:ch16-11} 叙述并证明星形区域上的 Poincaré 引理。
\item\label{exercise:ch16-12} 设 $(M,g)$ 是一个带边 Riemann 流形，$X$ 是一个光滑向量场。证明 Green 第一恒等式
\begin{align*}
\int_M(\operatorname{div}X)u\,\mathrm{d}V_g=\int_M\langle\operatorname{grad}u,X\rangle_g\,\mathrm{d}V_g-\int_{\partial M}u\langle X,N\rangle\,\mathrm{d}V_{\widetilde g}.
\end{align*}
并解释这与分部积分有何关系。

\item\label{exercise:ch16-13} 设 $(M,g)$ 是一个带边或不带边的 Riemann $n$ 流形。由 $\Delta u=-\operatorname{div}(\operatorname{grad}u)$ 定义的线性算子 $\Delta:C^\infty(M)\to C^\infty(M)$ 称为（几何）Laplace 算子。证明在任意光滑局部坐标下
\begin{align*}
\Delta u=-\frac{1}{\sqrt{\det g}}\frac{\partial}{\partial x^i}\left(g^{ij}\sqrt{\det g}\frac{\partial u}{\partial x^j}\right).
\end{align*}
由此推出在 $\mathbb{R}^n$ 上用欧氏度量与标准坐标时，$\Delta u=-\sum_{i=1}^n\partial^2u/(\partial x^i)^2$。[注：关于 Riemann 流形上 Laplace 算子的符号约定并无统一意见，许多作者把 $\Delta$ 定义为我们所定义算子的相反数。此处定义的几何 Laplace 算子与 $\mathbb{R}^n$ 上传统的 Laplace 算子符号相反，但它有两个明显的优点：本算子具有非负特征值（见习题~\cref{exercise:ch16-15}），并且与在微分形式上定义的 Laplace--Beltrami 算子一致（见习题~\cref{exercise:ch17-02}、\cref{exercise:ch17-03}）。]

\item\label{exercise:ch16-14} 设 $(M,g)$ 是一个带边或不带边的 Riemann 流形。若 $u\in C^\infty(M)$ 满足 $\Delta u=0$，则称 $u$ 是调和的。
\begin{enumerate}[(a)]
\item 设 $M$ 紧，证明 Green 恒等式：
\begin{align*}
\int_M u\Delta v\,\mathrm{d}V_g&=\int_M\langle\operatorname{grad}u,\operatorname{grad}v\rangle_g\,\mathrm{d}V_g-\int_{\partial M}uNv\,\mathrm{d}V_{\widetilde g},\\
\int_M(u\Delta v-v\Delta u)\,\mathrm{d}V_g&=\int_{\partial M}(vNu-uNv)\,\mathrm{d}V_{\widetilde g}.
\end{align*}
\item 证明：若 $M$ 紧、连通且 $\partial M=\varnothing$，则 $M$ 上的调和函数只有常数。
\item 证明：若 $M$ 紧、连通、$\partial M\neq\varnothing$，且 $u,v$ 是 $M$ 上在 $\partial M$ 上取值相同的调和函数，则 $u=v$。
\end{enumerate}

\item\label{exercise:ch16-15} 设 $(M,g)$ 是一个紧连通无边 Riemann 流形，$\Delta$ 是其几何 Laplace 算子。若存在一个不恒为零的光滑实值函数 $u$ 使 $\Delta u=\lambda u$，则称实数 $\lambda$ 为 $\Delta$ 的一个特征值，$u$ 为相应的特征函数。
\begin{enumerate}[(a)]
\item 证明 $0$ 是 $\Delta$ 的一个特征值，且其余特征值都严格为正。
\item 证明：若 $u,v$ 是对应于不同特征值的特征函数，则 $\int_M uv\,\mathrm{d}V_g=0$。
\end{enumerate}

\item\label{exercise:ch16-16} 设 $M$ 是一个具有非空边界的紧连通 Riemann $n$ 流形。若 $\Delta u=\lambda u$ 且 $u|_{\partial M}=0$，则称 $\lambda$ 为 Dirichlet 特征值；若 $\Delta u=\lambda u$ 且 $Nu|_{\partial M}=0$，则称 $\lambda$ 为 Neumann 特征值。
\begin{enumerate}[(a)]
\item 证明每个 Dirichlet 特征值都严格为正。
\item 证明 $0$ 是一个 Neumann 特征值，其余 Neumann 特征值都严格为正。
\end{enumerate}

\item\label{exercise:ch16-17} \textbf{Dirichlet 原理：}设 $M$ 是一个具有非空边界的紧连通 Riemann $n$ 流形。证明：$u\in C^\infty(M)$ 调和当且仅当它在所有边界值相同的光滑函数中使 $\int_M|\operatorname{grad}u|_g^2\,\mathrm{d}V_g$ 取最小。[提示：对任意在 $\partial M$ 上为零的函数 $f$，展开 $\int_M|\operatorname{grad}(u+\varepsilon f)|_g^2\,\mathrm{d}V_g$ 并利用习题~\cref{exercise:ch16-12}。]

\item\label{exercise:ch16-18} 设 $(M,g)$ 是一个定向 Riemann $n$ 流形。本题概述本章所定义算子 $\star$ 的一个重要推广。
\begin{enumerate}[(a)]
\item 对每个 $k=1,\dots,n$，证明 $g$ 在 $\Lambda^k(T_p^*M)$ 上确定唯一的内积 $\langle\cdot,\cdot\rangle_g$，满足
\begin{align*}
\langle\omega_1\wedge\cdots\wedge\omega_k,\eta_1\wedge\cdots\wedge\eta_k\rangle_g=\det\bigl(\langle(\omega_i)^\sharp,(\eta_j)^\sharp\rangle_g\bigr).
\end{align*}
\item 证明 Riemann 体积形式 $\mathrm{d}V_g$ 是关于此内积唯一的、定向为正的、范数为 $1$ 的 $n$-形式。
\item 对每个 $k=0,\dots,n$，证明存在唯一的光滑丛同态 $\star:\Lambda^kT^*M\to\Lambda^{n-k}T^*M$，满足 $\omega\wedge\star\eta=\langle\omega,\eta\rangle_g\,\mathrm{d}V_g$。这个映射称为 Hodge 星算子。
\item 证明 $\star f=f\,\mathrm{d}V_g$。
\item 证明 $\star\star\omega=(-1)^{k(n-k)}\omega$。
\end{enumerate}

\item\label{exercise:ch16-19} 把 $\mathbb{R}^n$ 看作具有欧氏度量与标准定向的 Riemann 流形。
\begin{enumerate}[(a)]
\item 计算 $\star\mathrm{d}x^i$。
\item 在 $n=4$ 时计算 $\star(\mathrm{d}x^i\wedge\mathrm{d}x^j)$。
\end{enumerate}

\item\label{exercise:ch16-20} 设 $M$ 是一个定向 Riemann $4$ 流形。若 $\star\omega=\omega$ 则称 $2$-形式 $\omega$ 是自对偶的；若 $\star\omega=-\omega$ 则称反自对偶的。
\begin{enumerate}[(a)]
\item 证明每个 $2$-形式都可唯一写成一个自对偶形式与一个反自对偶形式之和。
\item 在 $M=\mathbb{R}^4$ 上确定标准坐标中的自对偶与反自对偶形式。
\end{enumerate}

\item\label{exercise:ch16-21} 设 $(M,g)$ 是一个定向 Riemann 流形，$X\in\mathfrak X(M)$。证明
\begin{align*}
X\lrcorner\mathrm{d}V_g=\star X^\flat,\qquad \operatorname{div}X=\star\mathrm{d}\star X^\flat,
\end{align*}
且当 $\dim M=3$ 时 $\operatorname{curl}X=(\star\mathrm{d}X^\flat)^\sharp$。

\item\label{exercise:ch16-22} 设 $(M,g)$ 是一个紧定向 Riemann $n$ 流形。对 $1\leqslant k\leqslant n$，由 $d^*\omega=(-1)^{n(k+1)+1}\star\mathrm{d}\star\omega$ 定义 $d^*:\Omega^k(M)\to\Omega^{k-1}(M)$，并对 $0$-形式令 $d^*\omega=0$。
\begin{enumerate}[(a)]
\item 证明 $d^*\circ d^*=0$。
\item 证明 $(\omega,\eta)=\int_M\langle\omega,\eta\rangle_g\,\mathrm{d}V_g$ 定义了 $\Omega^k(M)$ 上的一个内积。
\item 证明 $(d^*\omega,\eta)=(\omega,d\eta)$。
\end{enumerate}

\item\label{exercise:ch16-23} 本题说明另一种证明某些 Riemann 度量不平坦的方法。令 $\mathbb{B}^2$ 为 $\mathbb{R}^2$ 中的单位圆盘，$g$ 为
\begin{align*}
g=\frac{\mathrm{d}x^2+\mathrm{d}y^2}{1-x^2-y^2}.
\end{align*}
\begin{enumerate}[(a)]
\item 证明：若 $g$ 平坦，则对充分小的 $r>0$，$g$-度量球 $\mathbb B_r^g(0)$ 的体积满足 $\operatorname{Vol}_g(\mathbb B_r^g(0))=\pi r^2$。
\item 对任意 $v\in\mathbb{B}^2$，通过计算从 $0$ 到 $v$ 的直线的 $g$-长度，证明 $d_g(0,v)\leqslant\tanh^{-1}|v|$，从而 $\mathbb B_r^g(0)$ 包含欧氏球 $\mathbb B_{\tanh r}(0)$。
\item 证明 $\operatorname{Vol}_g(\mathbb B_r^g(0))\geqslant\pi\sinh^2 r>\pi r^2$，从而 $g$ 不平坦。
\end{enumerate}
\end{enumerate}
\end{document}
